% Track A3 — multi-channel consistency of the matter-bounce f_NL
% SKELETON (2026-09-02). Results sections carry the real, committed numbers;
% Secs. II.C and V.B are stubs pending the compaction-function redo (A3-1).
% Every number here is traceable to research/track_a3_multichannel/outputs/*.json
% or to a cited abstract. Nothing is placeholder-numeric.
\documentclass[aps,prd,twocolumn,superscriptaddress,nofootinbib,longbibliography]{revtex4-2}

\usepackage{amsmath,amssymb}
\usepackage{graphicx}
\usepackage{booktabs}
\usepackage[colorlinks=true,linkcolor=blue,citecolor=blue,urlcolor=blue]{hyperref}

\newcommand{\fnl}{f_{\rm NL}}
\newcommand{\fnlloc}{f_{\rm NL}^{\rm local}}
\newcommand{\OmGW}{\Omega_{\rm GW}}
\newcommand{\fpbh}{f_{\rm PBH}}
\newcommand{\zetac}{\zeta_{\rm c}}

\begin{document}

\title{Multi-channel consistency of the matter-bounce non-Gaussianity:\\
pulsar-timing spectral slope, primordial black holes, and the
SPHEREx-era bispectrum at $\fnlloc = -35/16$}

\author{Houston Golden}
\email{houston@hubify.com}
\affiliation{Hubify Research, Big Bounce Cosmology Program}

\date{September 2, 2026}

\begin{abstract}
A dust-dominated contracting phase preceding a nonsingular bounce predicts a
parameter-free local non-Gaussianity of the primordial curvature perturbation.
The printed value of Cai \textit{et al.} is $\fnlloc = -35/8$; a
permutation-convention audit following Li \textit{et al.} gives exactly half of
it, $\fnlloc = -35/16 = -2.1875$, and that factor of two is not yet settled by an
independent derivation. We ask what the halved value does to three observational
channels, and we state each channel's result at exactly the strength its evidence
supports. (i)~\emph{Pulsar timing.} Refitting the NANOGrav 15-yr
Hellings--Downs-correlated Ceffyl free-spectrum posteriors, we obtain a
characteristic-strain spectral index $\gamma = 2.567 \pm 0.382$, consistent at
$1.14\sigma$ with the matter-bounce scalar-induced gravitational-wave prediction
$\gamma = 3$ ($\OmGW \propto f^2$), and disfavouring the supermassive-black-hole
binary value $\gamma = 13/3$ by a Savage--Dickey factor
$\log_{10} B_{\rm MB/SMBHB} = +3.85$. We emphasise that $\OmGW \propto f^2$ is
the causality-limited infrared slope of \emph{any} scalar-induced background, so
this channel is a consistency check and not a discriminator, and that it is
insensitive to the value of $\fnl$. A scale-invariant \emph{primordial} tensor
background ($n_{\rm T}=0$, $\gamma=5$) is excluded at $6.37\sigma$ by the same
posterior. (ii)~\emph{Primordial black holes.} Within a Press--Schechter
treatment with the local quadratic map, negative $\fnl$ imposes an absolute
ceiling $\zeta_{\max} \simeq -5/(12\fnl)$ on the curvature perturbation, which
therefore \emph{doubles} from $0.0952$ at $-35/8$ to $0.1905$ at $-35/16$. At an
amplitude calibrated so that the Gaussian case saturates $\fpbh = 1$, the halved
value yields $\fpbh = 7.3\times10^{-3}$ against $3.8\times10^{-6}$ at $-35/8$ --
a suppression weaker by three orders of magnitude, placing $-35/16$ inside rather
than far below the $10^{-3} \le \fpbh \le 1$ band reported for negative-$\fnl$
bounce scenarios. (iii)~\emph{Large-scale structure.} The prediction is
consistent with the DESI DR1 constraint $\fnlloc = -3.6^{+9.0}_{-9.1}$ at
$0.16\sigma$, and is the only channel sensitive to the value of $\fnl$: SPHEREx
reaches $3.13\sigma$ ($2.63\sigma$ after the bounce-to-local shape projection
$r = 0.84$) at its bispectrum-only forecast $\sigma_{\fnl} = 0.7$, and
$4.38\sigma$ ($3.68\sigma$) at its combined target $\sigma_{\fnl} = 0.5$. The
three channels are mutually consistent. The operative uncertainty is internal
rather than observational: only the bispectrum channel depends on which of
$-35/16$ or $-35/8$ is correct, and it depends on it strongly.
\end{abstract}

\maketitle

%======================================================================
\section{Introduction}
\label{sec:intro}

% STUB. To be written. Content to cover:
%  - the bounce-vs-inflation question and why a parameter-free f_NL matters;
%  - the matter-contraction f_NL and the -35/8 vs -35/16 factor of two;
%  - why a single number measured in three unrelated channels is the strongest
%    available near-term test;
%  - explicit statement that this paper reports a consistency assessment, not a
%    detection, and that one channel (PBH) is reported at Press-Schechter level
%    only.

The matter-bounce scenario -- a dust-dominated contracting phase matched through
a nonsingular bounce onto the standard expanding history -- predicts a
scale-invariant curvature power spectrum together with a definite, parameter-free
local bispectrum amplitude~\cite{Cai2009,CaiEasson2012,Quintin2015}. Section
\ref{sec:parameter} fixes the value we test; Secs.~\ref{sec:pta}--\ref{sec:lss}
take it through three channels; Sec.~\ref{sec:discussion} states what is and is
not established.

%======================================================================
\section{The parameter under test}
\label{sec:parameter}

Cai \textit{et al.}~\cite{Cai2009} print a squeezed-limit amplitude
$\fnlloc = -35/8 = -4.375$ for the matter bounce. Li \textit{et
al.}~\cite{Li2016}, whose no-go analysis extends that of Quintin \textit{et
al.}~\cite{Quintin2015}, obtain exactly one half of it,
\begin{equation}
  \fnlloc = -\frac{35}{16} = -2.1875 .
  \label{eq:fnl}
\end{equation}
The discrepancy traces to permutation counting in the in-in commutator and to a
$+(99/128)\sum_i k_i^3$ term carried by the Cai shape polynomial, whose removal
halves the bare squeezed limit exactly. We adopt Eq.~\eqref{eq:fnl} throughout
and carry $-35/8$ in parallel wherever a channel distinguishes them.

\textit{Scope statement.}---The factor of two has \emph{not} been settled. An
independent separate-universe (non-attractor $\delta N$) derivation performed for
this program returns $\fnlloc(\epsilon) = (5\epsilon-35)/8$, i.e. $-55/16$ at
$\epsilon = 3/2$, which is not a rational multiple of either published candidate
and therefore does not adjudicate between them; it does establish that the two
published values share a convention, so the discrepancy is not definitional. We
therefore never state a channel result at $-35/16$ alone where the two values
differ materially.

%======================================================================
\section{Channel I: the pulsar-timing spectral slope}
\label{sec:pta}

\subsection{Spectral conventions}

With the characteristic strain parameterised as $h_c(f) \propto
f^{(3-\gamma)/2}$ and
$\OmGW = (2\pi^2/3H_0^2)\, f^2 h_c^2(f)$, the energy density scales as
\begin{equation}
  \OmGW(f) \propto f^{\,5-\gamma}.
  \label{eq:slope}
\end{equation}
Hence $\gamma = 13/3$ for gravitational-wave-driven supermassive-black-hole
binary inspirals ($\OmGW \propto f^{2/3}$), $\gamma = 5$ for a scale-invariant
primordial tensor background ($n_{\rm T} = 0$, $\OmGW \propto f^0$), and
$\gamma = 3$ for $\OmGW \propto f^2$, the universal infrared scaling of the
scalar-induced background of a nonsingular matter bounce~\cite{Papanikolaou2025}.
We stress that $\gamma = 3$ is \emph{not} the $n_{\rm T} = 0$ case; the two are
separated by two units of slope.

\subsection{Posterior and Bayes factors}

We fit the NANOGrav 15-yr Hellings--Downs-correlated free-spectrum Ceffyl kernel
density posteriors~\cite{NANOGrav15,NANOGravKDE} ($30$ frequency bins,
$T_{\rm obs} = 16.03\,$yr) with a power law $(\gamma, \log_{10}A)$ under priors
$\gamma \sim U[0,7]$, $\log_{10}A \sim U[-18,-11]$, using $32$ walkers and
$10^4$ production steps after $2500$ burn-in ($3.2\times10^5$ samples,
${\rm ESS} = 5.5\times10^3$). The marginal is
\begin{equation}
  \gamma = 2.567 \pm 0.382, \qquad
  \log_{10} A = -14.03 \pm 0.38 ,
  \label{eq:gamma}
\end{equation}
with median $2.591$ and $68\%$ interval $[2.304, 2.882]$.

Savage--Dickey density ratios for the nested point hypotheses $\gamma = \gamma_*$
against free $\gamma$, $B_{\gamma_*/{\rm free}} = p(\gamma_*|d)/\pi(\gamma_*)$
with $\pi = 1/7$, are collected in Table~\ref{tab:pta}.

\begin{table}[t]
\caption{\label{tab:pta}Channel I. Savage--Dickey factors and Gaussian
$z$-distances from the posterior of Eq.~\eqref{eq:gamma}. All values reproduced
from the committed chain
(\texttt{pta\_gamma\_reproduce.py}); differences from the archived record are
$\le 3\times10^{-15}$.}
\begin{ruledtabular}
\footnotesize
\begin{tabular}{lcccc}
Hypothesis & $\gamma_*$ & $\OmGW$ slope & $B_{\gamma_*/{\rm free}}$ & $z$ \\
\colrule
Matter bounce (induced) & $3$      & $f^{2}$   & $3.23$                & $1.14\sigma$ \\
SMBHB inspiral          & $13/3$   & $f^{2/3}$ & $4.52\times10^{-4}$   & $4.63\sigma$ \\
Primordial tensors      & $5$      & $f^{0}$   & $1.86\times10^{-24}$  & $6.37\sigma$ \\
\end{tabular}
\end{ruledtabular}
\end{table}

The matter-bounce-to-SMBHB ratio is
$B_{\rm MB/SMBHB} = 7.14\times10^{3}$, i.e. $\log_{10}B = +3.85$. Narrowing the
$\gamma$ prior from $U[0,7]$ to $U[2,5]$ moves $B_{3/{\rm free}}$ from $3.23$ to
$1.47$ while leaving $B_{\rm MB/SMBHB}$ stable in the range
$7.1$--$8.7\times10^3$; the nested factor is prior-sensitive, the model ratio is
not.

\subsection{What this channel does and does not establish}

The measured slope is consistent with the matter-bounce induced-GW prediction and
strongly disfavours both the SMBHB expectation and a scale-invariant primordial
tensor background. It does not favour the matter bounce \emph{specifically}:
$\OmGW \propto f^2$ is the causality-limited infrared slope common to
scalar-induced backgrounds of essentially any origin. Nor does this channel
constrain $\fnl$: it measures a slope, not a bispectrum amplitude, and is
identical at $-35/16$ and $-35/8$. Converting it into a discriminating test
requires propagating the matter-bounce scalar spectrum through the induced-GW
kernel to predict an \emph{amplitude} in the nHz band; that computation is not
attempted here.

%======================================================================
\section{Channel II: primordial black hole abundance}
\label{sec:pbh}

\subsection{Press--Schechter with local quadratic non-Gaussianity}

Writing $\zeta = \zeta_G + A(\zeta_G^2 - \sigma^2)$ with $A = \tfrac{3}{5}\fnl$
and $\zeta_G \sim \mathcal{N}(0,\sigma^2)$~\cite{YoungByrnes2013,Franciolini2018},
the mass fraction is $\beta(\zetac) = P(\zeta > \zetac)$, obtained exactly by
inverting the quadratic. Present-day abundance follows the standard radiation-era
relation~\cite{Sasaki2018}
\begin{equation}
\begin{split}
  \fpbh(M) = 1.68\times10^{8}\,
  &\Big(\frac{\gamma_{\rm c}}{0.2}\Big)^{1/2}
  \Big(\frac{g_*}{106.75}\Big)^{-1/4}\\
  &\times\Big(\frac{M}{M_\odot}\Big)^{-1/2}\beta(M),
\end{split}
\end{equation}
with $\gamma_{\rm c} = 0.2$, $g_* = 106.75$, and $M = 10^{20}\,$g.

\subsection{An absolute ceiling that doubles}

For $\fnl < 0$ the map is a downward parabola, so
\begin{equation}
  \zeta_{\max} = -\frac{5}{12 \fnl} + \frac{3}{5}|\fnl|\sigma^2 ,
  \label{eq:ceiling}
\end{equation}
and no realisation exceeds it. The leading term scales as $1/|\fnl|$:
$\zeta_{\max} = 0.09524$ at $-35/8$ and $0.19048$ at $-35/16$, a ratio of exactly
$2$. This is the sharpest statement the channel makes about the two candidate
values.

\subsection{Abundance at fixed amplitude}

The curvature power-spectrum amplitude is set by the source, and $\fnl$ then
decides whether black holes are overproduced. Calibrating $\sigma$ on the
Gaussian case so that $\fpbh = 1$ and reading off the two matter-bounce values at
that same $\sigma$ gives Table~\ref{tab:pbh}.

\begin{table}[t]
\caption{\label{tab:pbh}Channel II. $\fpbh$ at $M = 10^{20}\,$g, with $\sigma$
calibrated at each threshold so that the Gaussian case gives $\fpbh = 1$.
Entries ``$0$'' mean $\zetac$ exceeds the ceiling of Eq.~\eqref{eq:ceiling}.
Perturbativity $0.6|\fnl|\sigma \le 0.04$ throughout.}
\begin{ruledtabular}
\begin{tabular}{ccccc}
$\zetac$ & $\sigma_*$ & $\fpbh(-35/16)$ & $\fpbh(-35/8)$ & ratio \\
\colrule
$0.05$ & $0.006325$ & $7.32\times10^{-3}$ & $3.75\times10^{-6}$  & $1.95\times10^{3}$ \\
$0.08$ & $0.010120$ & $1.24\times10^{-4}$ & $1.12\times10^{-14}$ & $1.11\times10^{10}$ \\
$0.10$ & $0.012650$ & $3.75\times10^{-6}$ & $0$ & --- \\
$0.12$ & $0.015179$ & $4.29\times10^{-8}$ & $0$ & --- \\
$0.15$ & $0.018974$ & $1.82\times10^{-12}$ & $0$ & --- \\
\end{tabular}
\end{ruledtabular}
\end{table}

Negative $\fnl$ suppresses black hole formation, and $-35/16$ is a substantially
weaker suppressor than $-35/8$: three to ten orders of magnitude at fixed
amplitude. At $\zetac = 0.05$ the halved value lands inside the
$10^{-3}\le\fpbh\le1$ band reported for negative-$\fnl$ bounce
scenarios~\cite{Choudhury2025}, while $-35/8$ falls three orders below it.

\subsection{Limitation of the quadratic truncation}
\label{sec:pbh_limit}

% STUB (item A3-1). To be replaced by a compaction-function calculation.
At the standard curvature thresholds $\zetac \approx 0.45$--$1$, the ceiling of
Eq.~\eqref{eq:ceiling} forces $\beta = 0$ identically for \emph{both} negative
values in the rare-tail regime. This is a property of the truncated quadratic
map, not a physical statement, and it is why the compaction-function formation
criterion is used in the literature~\cite{Choudhury2025}. Choudhury \textit{et
al.} treat $\fnl = (-39.95, -35/8)$ in an effective-field-theory bounce with an
ultra-slow-roll phase and report $10^{-3}\le\fpbh\le1$, complete mitigation of
overproduction, and a perturbativity bound $|\fnl| \lesssim 60$; they do not
treat $-35/16$. Redoing that calculation at $-35/16$ is the principal open item
of this work and is \emph{not} claimed here.

%======================================================================
\section{Channel III: the large-scale-structure bispectrum}
\label{sec:lss}

\subsection{Current constraints}

The DESI DR1 analysis of $1{,}631{,}716$ luminous red galaxies
($0.6<z<1.1$) and $1{,}189{,}129$ quasars
($0.8<z<3.1$)~\cite{Chaussidon2024} gives
$\fnlloc = -3.6^{+9.0}_{-9.1}$ at $68\%$ confidence with a merger model for the
quasar response, or $\fnlloc = +3.5^{+10.7}_{-7.4}$ assuming universality. The
prediction of Eq.~\eqref{eq:fnl} sits $0.16\sigma$ from the former and
$0.77\sigma$ from the latter. The discriminating power is
$|\fnl|/\sigma = 0.24$: DESI DR1 is consistent with the prediction and cannot
test it.

\subsection{Forecast reach}

The matter-bounce bispectrum is not exactly the local shape; we adopt the
noise-weighted overlap $r = 0.84$, so an experiment quoting $\sigma_{\rm local}$
constrains the bounce amplitude at $\sigma_{\rm local}/r$.
Table~\ref{tab:reach} gives both.

\begin{table}[t]
\caption{\label{tab:reach}Channel III. Detection significance of
$\fnlloc = -2.1875$. Forecast $\sigma$ values are quoted from the cited
abstracts; the significances are $|\fnl|/\sigma$ and $|\fnl|r/\sigma$.
MegaMapper is an unapproved design and is quoted only at its published
order-unity level; no systematic budget is transferred to it.}
\begin{ruledtabular}
\begin{tabular}{lccc}
Survey & $\sigma_{\fnl}$ & bare & $r$-projected \\
\colrule
SPHEREx, bispectrum only~\cite{Heinrich2023} & $0.7$ & $3.13\sigma$ & $2.63\sigma$ \\
SPHEREx, $P+B$ target~\cite{Heinrich2023,Dore2014} & $0.5$ & $4.38\sigma$ & $3.68\sigma$ \\
MegaMapper class~\cite{Ferraro2019,Sailer2021} & $\sim 1$ & $2.19\sigma$ & $1.84\sigma$ \\
\end{tabular}
\end{ruledtabular}
\end{table}

At $-35/8$ the same rows read $6.25\sigma$, $8.75\sigma$ and $4.38\sigma$ bare.
This is the only channel in which the factor of two is observationally
consequential.

%======================================================================
\section{Discussion}
\label{sec:discussion}

\subsection{The joint statement}

The three channels are mutually consistent at $\fnlloc = -35/16$. The pulsar
timing slope agrees at $1.14\sigma$ and rejects the two competing spectral shapes;
the primordial black hole channel is suppressed but not extinguished, and lands
in the sizeable-abundance band; the bispectrum prediction is consistent with the
best current measurement and reachable by SPHEREx. No channel is in tension with
another.

\subsection{Where the consistency is weak}

Two of the three channels are insensitive to the parameter. The pulsar-timing
result is a statement about a slope that any scalar-induced background shares,
and it is identical at both candidate $\fnl$ values. The black hole channel is
reported only at Press--Schechter level. The operative uncertainty is therefore
internal: only the bispectrum channel separates $-35/16$ from $-35/8$, and it
separates them by more than $3\sigma$ at SPHEREx sensitivity, so settling the
factor of two by an independent derivation is a prerequisite for a sharp
prediction rather than a bookkeeping detail.

\subsection{Next steps}
\label{sec:next}

% STUB. To be completed alongside items A3-1 through A3-5.
(i)~a compaction-function recomputation of Channel II at $-35/16$;
(ii)~an independent gradient-expansion or $\delta N$ derivation of
Eq.~\eqref{eq:fnl};
(iii)~propagation of the matter-bounce scalar spectrum through the induced-GW
kernel to predict the nHz \emph{amplitude}, converting Channel I from a
consistency check into a test;
(iv)~re-derivation of the shape overlap $r$ at the $-35/16$ fiducial.

%======================================================================
\section*{Data and code availability}
\sloppy

All three channel computations are committed at
\texttt{research/track\_a3\_multichannel/}: \texttt{pta\_gamma\_reproduce.py},
\texttt{pbh\_abundance\_fnl.py}, \texttt{survey\_reach\_fnl.py}, with results in
\texttt{outputs/} and reproducibility manifests at
\texttt{reproducibility/manifests/}\allowbreak\texttt{experiments/a3-*.json}.
The pulsar-timing chain is
\texttt{pipelines/p3\_pta\_mcmc/}\allowbreak
\texttt{free\_spectrum\_real\_2026-05-01/}\allowbreak
\texttt{chain\_real\_freespec.npy}
(SHA-256 \texttt{50abc38a\ldots10fc}), derived from the NANOGrav 15-yr KDE free
spectra, Zenodo \texttt{10.5281/zenodo.8060824}. Every computation ran on a local
CPU in under $0.1\,$s at zero compute cost.

%======================================================================
\begin{thebibliography}{99}

\bibitem{Cai2009} Y.-F. Cai, W. Xue, R. Brandenberger, and X. Zhang,
  ``Non-Gaussianity in a Matter Bounce,'' JCAP \textbf{0905}, 011 (2009),
  arXiv:0903.0631.

\bibitem{CaiEasson2012} Y.-F. Cai, D. A. Easson, and R. Brandenberger,
  ``Towards a Nonsingular Bouncing Cosmology,'' JCAP \textbf{1208}, 020 (2012),
  arXiv:1206.2382.

\bibitem{Quintin2015} J. Quintin, Z. Sherkatghanad, Y.-F. Cai, and
  R. H. Brandenberger, ``Evolution of cosmological perturbations and the
  production of non-Gaussianities through a nonsingular bounce: Indications for a no-go theorem in single field matter bounce cosmologies,''
  Phys. Rev. D \textbf{92}, 063532 (2015), arXiv:1508.04141.

\bibitem{Li2016} Y.-B. Li, J. Quintin, D.-G. Wang, and Y.-F. Cai,
  ``Matter bounce cosmology with a generalized single field: non-Gaussianity and
  an extended no-go theorem,'' JCAP \textbf{1703}, 031 (2017),
  arXiv:1612.02036.

\bibitem{Papanikolaou2025} T. Papanikolaou,
  ``Gravitational wave signatures of non-singular matter bouncing cosmology in
  NANOGrav and beyond,'' arXiv:2504.11641.

\bibitem{NANOGrav15} G. Agazie \textit{et al.} (NANOGrav Collaboration),
  ``The NANOGrav 15 yr Data Set: Evidence for a Gravitational-wave Background,''
  Astrophys. J. Lett. \textbf{951}, L8 (2023), arXiv:2306.16213.

\bibitem{NANOGravKDE} NANOGrav Collaboration,
  ``NANOGrav 15 yr KDE Free Spectra v1.0.0,''
  Zenodo, \href{https://doi.org/10.5281/zenodo.8060824}{10.5281/zenodo.8060824}.

\bibitem{YoungByrnes2013} S. Young and C. T. Byrnes,
  ``Primordial black holes in non-Gaussian regimes,'' JCAP \textbf{1308}, 052
  (2013), arXiv:1307.4995.

\bibitem{Franciolini2018} G. Franciolini, A. Kehagias, S. Matarrese, and
  A. Riotto, ``Primordial Black Holes from Inflation and non-Gaussianity,''
  JCAP \textbf{1803}, 016 (2018), arXiv:1801.09415.

\bibitem{Sasaki2018} M. Sasaki, T. Suyama, T. Tanaka, and S. Yokoyama,
  ``Primordial black holes -- perspectives in gravitational wave astronomy,''
  Class. Quant. Grav. \textbf{35}, 063001 (2018), arXiv:1801.05235.

\bibitem{Choudhury2025} S. Choudhury, K. Dey, S. Ganguly, A. Karde,
  S. K. Singh, and P. Tiwari,
  ``Negative non-Gaussianity as a salvager for PBHs with PTAs in bounce,''
  Eur. Phys. J. C \textbf{85}, 472 (2025), arXiv:2409.18983.

\bibitem{Chaussidon2024} E. Chaussidon \textit{et al.} (DESI Collaboration),
  ``Constraining primordial non-Gaussianity with DESI 2024 LRG and QSO
  samples,'' arXiv:2411.17623.

\bibitem{Heinrich2023} C. Heinrich, O. Dor\'e, and E. Krause,
  ``Measuring $f_{\rm NL}$ with the SPHEREx Multi-tracer Redshift Space
  Bispectrum,'' arXiv:2311.13082.

\bibitem{Dore2014} O. Dor\'e \textit{et al.},
  ``Cosmology with the SPHEREx All-Sky Spectral Survey,'' arXiv:1412.4872.

\bibitem{Ferraro2019} S. Ferraro \textit{et al.},
  ``Inflation and Dark Energy from spectroscopy at $z>2$,''
  Astro2020 science white paper, arXiv:1903.09208.

\bibitem{Sailer2021} N. Sailer, E. Castorina, S. Ferraro, and M. White,
  ``Cosmology at high redshift -- a probe of fundamental physics,''
  JCAP \textbf{12}, 049 (2021), arXiv:2106.09713.

\end{thebibliography}

\end{document}
